%+TITLE: Modelo de Weinberg--Salam
%+DESCRIPTION: Construimos la teoría electrodébil, las masas y los acoplamientos de Yukawa, y reunimos el lagrangiano del modelo estándar
% Fuente: transcripcion_teoria_clasica_campos_paginas_36-43/transcripcion_paginas_36-43.tex
% Páginas manuscritas 40--43.
% Revisión editorial: quiralidad, cargas, potencial de Higgs y acción final.
\documentclass[11pt,a4paper]{article}

\usepackage[utf8]{inputenc}
\usepackage[T1]{fontenc}
\usepackage[spanish,es-nodecimaldot]{babel}
\usepackage{lmodern}
\usepackage{amsmath,amssymb,mathtools}
\usepackage{graphicx}
\usepackage{xcolor}
\usepackage{geometry}
\usepackage{enumitem}
\usepackage{microtype}

\geometry{margin=2.3cm}
\setlength{\parindent}{0pt}
\setlength{\parskip}{0.65em}
\setlist[itemize]{leftmargin=1.6em,itemsep=0.25em,topsep=0.25em}

\definecolor{title-red}{RGB}{165,25,20}
\definecolor{concept-green}{RGB}{20,125,45}
\newcommand{\concept}[1]{\textcolor{concept-green}{#1}}

\newcommand{\ket}[1]{\left|#1\right\rangle}
\newcommand{\bra}[1]{\left\langle#1\right|}
\newcommand{\braket}[2]{\left\langle #1\middle|#2\right\rangle}
\newcommand{\hc}{\mathrm{h.c.}}
\newcommand{\Tr}{\operatorname{Tr}}
\newcommand{\Diffs}{\mathrm{Diffs}}
\newcommand{\slashedpartial}{\not\!\partial}

\begin{document}

% -----------------------------------------------------------------------------
% Página manuscrita 40
% -----------------------------------------------------------------------------
Ahora describiremos el modelo que describe las interacciones nucleares débiles.
En realidad, según este modelo, estas interacciones están unificadas con la
electrodinámica en lo que se llama teoría electrodébil.

Para empezar, definimos la matriz
\[
  \gamma^5=i\gamma^0\gamma^1\gamma^2\gamma^3
  =\begin{pmatrix}
    -\mathbb 1_{2\times2}&0\\
    0&\mathbb 1_{2\times2}
  \end{pmatrix}.
\]
De su definición o forma explícita vemos que
\[
  \{\gamma^5,\gamma^\mu\}=0,\qquad(\gamma^5)^2=\mathbb1.
\]

Podemos usar esta matriz para definir los proyectores
\[
  P_R=\frac12(1+\gamma^5)
  =\begin{pmatrix}0&0\\0&\mathbb1_{2\times2}\end{pmatrix},
  \qquad
  P_L=\frac12(1-\gamma^5)
  =\begin{pmatrix}\mathbb1_{2\times2}&0\\0&0\end{pmatrix}.
\]

Es fácil ver que
\[
  P_R^\dagger=P_R,
  \qquad
  P_L^\dagger=P_L,
  \qquad
  P_R^2=P_R,
  \qquad
  P_L^2=P_L,
  \qquad
  P_RP_L=0,
  \qquad
  P_R+P_L=\mathbb 1_{4\times4},
\]
\[
  P_R\gamma^\mu=\gamma^\mu P_L.
\]

Aplicados sobre los espinores de Dirac, estos proyectores proyectan la parte
izquierda o derecha:
\[
  P_R\psi=P_R\begin{pmatrix}\chi\\\xi\end{pmatrix}
  =\begin{pmatrix}0\\\xi\end{pmatrix}\equiv\psi_R,
  \qquad
  P_L\psi=P_L\begin{pmatrix}\chi\\\xi\end{pmatrix}
  =\begin{pmatrix}\chi\\0\end{pmatrix}\equiv\psi_L.
\]

Note que
\begin{align*}
  P_R\gamma^0\gamma^\mu
  &=\frac12(1+\gamma^5)\gamma^0\gamma^\mu
   =\frac12\gamma^0(1-\gamma^5)\gamma^\mu\\
  &=\frac12\gamma^0\gamma^\mu(1+\gamma^5)
   =\gamma^0\gamma^\mu P_R,
\end{align*}
\[
  P_L\gamma^0\gamma^\mu=\gamma^0\gamma^\mu P_L.
\]

Las barras de los campos quirales se definen después de proyectar:
$\bar\psi_L=\overline{\psi_L}=\bar\psi P_R$ y
$\bar\psi_R=\overline{\psi_R}=\bar\psi P_L$.

Entonces el término cinético de Dirac se puede escribir
\begin{align*}
  i\bar\psi\slashedpartial\psi
  &=i\psi^\dagger\gamma^0\gamma^\mu\partial_\mu\psi\\
  &=i\psi^\dagger\gamma^0\gamma^\mu(P_R+P_L)\partial_\mu\psi\\
  &=i\psi^\dagger\gamma^0\gamma^\mu(P_R^2+P_L^2)\partial_\mu\psi\\
  &=i\psi^\dagger P_R\gamma^0\gamma^\mu P_R\partial_\mu\psi
    +i\psi^\dagger P_L\gamma^0\gamma^\mu P_L\partial_\mu\psi\\
  &=i\psi_R^\dagger\gamma^0\gamma^\mu\partial_\mu\psi_R
    +i\psi_L^\dagger\gamma^0\gamma^\mu\partial_\mu\psi_L\\
  &=i\bar\psi_L\slashedpartial\psi_L
    +i\bar\psi_R\slashedpartial\psi_R.
\end{align*}

Ahora estudiemos una teoría con electrones, descritos con el espinor de Dirac
$e$, y neutrinos electrónicos, descritos por un espinor de Weyl izquierdo
escrito en notación de cuatro componentes, $\nu=\nu_L$, con
\[
  P_R\nu=0.
\]
Los tomamos por ahora sin masa y escribimos
\[
  \mathcal L_0
  =i\bar e_L\slashedpartial e_L
   +i\bar e_R\slashedpartial e_R
   +i\bar\nu_L\slashedpartial\nu_L.
\]

Según el modelo de Weinberg--Salam, la interacción débil está unificada con la
electrodinámica y corresponde a una simetría $SU(2)\times U(1)$. Las matrices
$2\times2$ (de la representación fundamental) del grupo $SU(2)$ actúan sobre el
doblete de espinores izquierdos, pero no los derechos:
\[
  L=\begin{pmatrix}\nu_L\\e_L\end{pmatrix},
  \qquad
  L\longmapsto e^{ig\vec\alpha\cdot\vec\sigma/2}L,
  \qquad
  e_R\longmapsto e_R.
\]

En términos de $L$ el lagrangiano queda
\[
  \mathcal L=i\bar L\slashedpartial L+i\bar e_R\slashedpartial e_R.
\]
A veces a esta transformación se la llama de \concept{isospín débil}. Decimos
que $L$ tiene isospín $1/2$ (pertenece a esa representación de $SU(2)$) y
$e_R$ tiene isospín $0$. Vemos además que hay una relación entre la carga y la
tercera componente de isospín:
\[
  L:\qquad Q_{\mathrm{em}}=I_3-\frac12,
  \qquad
  e_R:\qquad Q_{\mathrm{em}}=I_3-1.
\]

% -----------------------------------------------------------------------------
% Página manuscrita 41
% -----------------------------------------------------------------------------
La simetría $U(1)$, llamada de \concept{hipercarga}, actúa de la siguiente
manera:
\[
  L\longmapsto e^{-ig'\beta/2}L,
  \qquad
  e_R\longmapsto e^{-ig'\beta}e_R.
\]

Decimos que las hipercargas son
\[
  Y_L=-\frac12,
  \qquad
  Y_{e_R}=-1.
\]
El número cuántico de carga eléctrica es
\[
  Q_{\mathrm{em}}=I_3+Y,\qquad q=e_{\mathrm{em}}Q_{\mathrm{em}},
  \qquad e_{\mathrm{em}}>0.
\]

Según el modelo la simetría es de gauge, tal que debemos introducir derivadas
covariantes y campos de gauge:
\[
  D_\mu=\partial_\mu-ig'YX_\mu-ig\vec J\cdot\vec W_\mu,
\]
\[
  D_\mu L
  =\partial_\mu L+\frac{ig'}2X_\mu L
   -\frac{ig}2\vec\sigma\cdot\vec W_\mu L,
  \qquad
  D_\mu e_R=\partial_\mu e_R+ig'X_\mu e_R.
\]

El lagrangiano con esta invariancia de gauge es
\[
  \mathcal L_0
  =i\bar L\not\!D L+i\bar e_R\not\!D e_R
   -\frac14X_{\mu\nu}X^{\mu\nu}
   -\frac14\vec W_{\mu\nu}\cdot\vec W^{\mu\nu},
\]
con
\[
  X_{\mu\nu}=\partial_\mu X_\nu-\partial_\nu X_\mu,
\]
\[
  W_{\mu\nu}^i
  =\partial_\mu W_\nu^i-\partial_\nu W_\mu^i
   +g\epsilon_{ijk}W_\mu^jW_\nu^k.
\]

Por ahora todos los campos son sin masa. De hecho, si intentamos escribir un
término de masa para el electrón $\bar e_Re_L$ no es invariante bajo la
simetría. Para dar masas a los bosones de gauge y al electrón introducimos un
campo de Higgs:
\[
  \mathcal L=\mathcal L_0+\mathcal L_{\mathrm{Higgs}},
\]
\[
  \mathcal L_{\mathrm{Higgs}}
  =-(D_\mu\phi)^\dagger D^\mu\phi
   -m^2\phi^\dagger\phi
   -\lambda(\phi^\dagger\phi)^2.
\]

El campo de Higgs lo tomamos en la representación fundamental de $SU(2)$ y con
$Y=1/2$:
\[
  \phi=\begin{pmatrix}\phi^+\\\phi^0\end{pmatrix}.
\]

Para romper la simetría tomamos $m^2<0$ y $\lambda>0$. El mínimo satisface
\[
  \phi^\dagger\phi=|\phi^+|^2+|\phi^0|^2
  =-\frac{m^2}{2\lambda}\equiv v^2.
\]
Nuestra $v$ es el valor de la componente neutra del doblete;
la convención $\phi_0=(0,v_{\mathrm{SM}}/\sqrt2)^T$ usa
$v_{\mathrm{SM}}=\sqrt2v$.
Sin pérdida de generalidad,
\[
  \phi_0=\begin{pmatrix}0\\v\end{pmatrix},
  \qquad
  \phi=
  \begin{pmatrix}
    \dfrac1{\sqrt2}(\phi_3+i\phi_4)\\[2mm]
    v+\dfrac1{\sqrt2}(\phi_1+i\phi_2)
  \end{pmatrix}.
\]

Aplicamos el mismo truco de escoger el gauge (unitario) tal que
$\phi_2=\phi_3=\phi_4=0$, tal que
\[
  \phi=\begin{pmatrix}0\\v+\dfrac1{\sqrt2}h\end{pmatrix}.
\]
Entonces
\[
  D_\mu\phi
  =\left[
    \partial_\mu-\frac{ig'}2X_\mu
    -\frac{ig}{2}
    \begin{pmatrix}
      W_\mu^3&W_\mu^1-iW_\mu^2\\
      W_\mu^1+iW_\mu^2&-W_\mu^3
    \end{pmatrix}
  \right]
  \begin{pmatrix}0\\v+\dfrac1{\sqrt2}h\end{pmatrix},
\]
por lo que
\[
  D_\mu\phi=
  \begin{pmatrix}
    -\dfrac{ig}{2}\left(v+\dfrac h{\sqrt2}\right)
      (W_\mu^1-iW_\mu^2)\\[3mm]
    \dfrac1{\sqrt2}\partial_\mu h
    -\dfrac{ig'}2X_\mu\left(v+\dfrac h{\sqrt2}\right)
    +\dfrac{ig}{2}W_\mu^3\left(v+\dfrac h{\sqrt2}\right)
  \end{pmatrix}.
\]

\begin{align*}
  (D_\mu\phi)^\dagger(D^\mu\phi)
  ={}&\frac{g^2}{4}\left(v+\frac h{\sqrt2}\right)^2
  (W_\mu^1+iW_\mu^2)(W^{1\mu}-iW^{2\mu})
  +\frac12\partial_\mu h\partial^\mu h\\
  &+\frac14\left(v+\frac h{\sqrt2}\right)^2
  (gW_\mu^3-g'X_\mu)(gW^{3\mu}-g'X^\mu)\\
  ={}&\frac{g^2v^2}{4}
  \left(W_\mu^1W^{1\mu}+W_\mu^2W^{2\mu}\right)
  +\frac12\partial_\mu h\partial^\mu h\\
  &+\frac{v^2}{4}
  (gW_\mu^3-g'X_\mu)(gW^{3\mu}-g'X^\mu)
  +\text{cúbicos y cuárticos}.
\end{align*}

Vemos que aparece un término cinético para $h$ junto con términos de masa para
$W_\mu^1$, $W_\mu^2$ y para la combinación $gW_\mu^3-g'X_\mu$.

% -----------------------------------------------------------------------------
% Página manuscrita 42
% -----------------------------------------------------------------------------
La combinación con masa recibe un nombre:
\[
  Z_\mu
  =\frac{gW_\mu^3-g'X_\mu}{(g'^2+g^2)^{1/2}}
  =\cos\theta_WW_\mu^3-\sin\theta_WX_\mu,
  \qquad
  \tan\theta_W=\frac{g'}g.
\]

Los bosones relacionados con la interacción débil son $W_\mu^1$, $W_\mu^2$ y
$Z_\mu$, con masas
\[
  m_W^2=\frac{g^2v^2}{2},
  \qquad
  m_Z^2=\frac{(g^2+g'^2)v^2}{2}.
\]

La combinación ortogonal a $Z_\mu$ no tiene masa y la identificamos con el
fotón:
\[
  A_\mu
  =\frac{g'W_\mu^3+gX_\mu}{(g^2+g'^2)^{1/2}}
  =\sin\theta_WW_\mu^3+\cos\theta_WX_\mu.
\]

Para hacer esa identificación notemos que podemos escribir
\[
  X_\mu=\cos\theta_WA_\mu-\sin\theta_WZ_\mu,
\]
\[
  W_\mu^3=\cos\theta_WZ_\mu+\sin\theta_WA_\mu.
\]

Volvamos a los términos para el electrón y el neutrino:
\begin{align*}
  D_\mu
  \begin{pmatrix}\nu_L\\e_L\end{pmatrix}
  ={}&\partial_\mu
  \begin{pmatrix}\nu_L\\e_L\end{pmatrix}
  +\frac{ig'}2\cos\theta_WA_\mu
  \begin{pmatrix}\nu_L\\e_L\end{pmatrix}
  -\frac{ig}2\sin\theta_WA_\mu
  \begin{pmatrix}1&0\\0&-1\end{pmatrix}
  \begin{pmatrix}\nu_L\\e_L\end{pmatrix}\\
  &+\text{términos con }Z_\mu,W_\mu^1,W_\mu^2\\
  ={}&
  \begin{pmatrix}
    \partial_\mu\nu_L\\[1mm]
    \partial_\mu e_L+
    i\dfrac{gg'}{(g^2+g'^2)^{1/2}}A_\mu e_L
  \end{pmatrix}
  +\cdots,
\end{align*}
\[
  D_\mu e_R
  =\partial_\mu e_R+ig'\cos\theta_WA_\mu e_R
   +\text{términos con }Z_\mu
  =\partial_\mu e_R+i\frac{gg'}{(g^2+g'^2)^{1/2}}A_\mu e_R+\cdots.
\]

El signo elegido para $A_\mu$ hace que la derivada electromagnética sea
$D_\mu=\partial_\mu-iqA_\mu+\cdots$, como en electrodinámica, con
\[
  e_{\mathrm{em}}=\frac{gg'}{\sqrt{g^2+g'^2}},
  \qquad q_e=-e_{\mathrm{em}},\qquad q_\nu=0.
\]
Ambas quiralidades del electrón tienen la misma carga.
Los bosones cargados son
$W^\pm_\mu=(W^1_\mu\mp iW^2_\mu)/\sqrt2$.
La expansión del potencial también da $m_h^2=4\lambda v^2$.

Podemos agregarle más términos al lagrangiano invariantes bajo la simetría y con
un coeficiente adimensional:
\[
  \mathcal L
  =\mathcal L_0+\mathcal L_{\mathrm{Higgs}}+\mathcal L_{\mathrm{Yukawa}},
  \qquad
  \mathcal L_{\mathrm{Yukawa}}
  =-\left(y_e\bar L\phi e_R+\hc\right).
\]

Después de la ruptura de simetría nos queda
\begin{align*}
  \mathcal L_{\mathrm{Yukawa}}
  &=-y_e(\bar\nu_L\ \bar e_L)
    \begin{pmatrix}0\\v+\dfrac1{\sqrt2}h\end{pmatrix}e_R
    +\text{h.c.}\\
  &=-y_ev\bar e_Le_R-\frac{y_eh}{\sqrt2}\bar e_Le_R+\text{h.c.}
\end{align*}

Tras una redefinición de fase podemos tomar $y_e$ real y positivo.
Entonces $m_e=y_ev$ y el término de masa es $-m_e\bar e e$.

Los quarks también tienen interacción electrodébil con unas diferencias. Los
quarks $u$ y $d$ izquierdos forman el doblete, y ambos tienen versión derecha.
Las hipercargas son
\[
  Q=\begin{pmatrix}u_L\\d_L\end{pmatrix},
  \qquad
  Y_Q=\frac16,
  \qquad
  Y_{u_R}=\frac23,
  \qquad
  Y_{d_R}=-\frac13.
\]

Para dar masa al quark de tipo up necesitamos el doblete conjugado
$\widetilde\phi=i\sigma_2\phi^*$, con hipercarga $-1/2$.
Los acoplamientos de Yukawa son
\[
  \mathcal L_{\mathrm{Yukawa,quarks}}
  =-\left(y_d\bar Q\phi d_R+y_u\bar Q\widetilde\phi u_R+\hc\right).
\]

% -----------------------------------------------------------------------------
% Página manuscrita 43
% -----------------------------------------------------------------------------
Después de la ruptura de simetría,
$\phi_0=(0,v)^T$ y $\widetilde\phi_0=(v,0)^T$, de modo que
\[
  \mathcal L_{\mathrm{Yukawa,quarks}}
  \supset-\left(y_dv\bar d_Ld_R+y_uv\bar u_Lu_R+\hc\right).
\]
Para una familia, las masas son $m_d=|y_d|v$ y $m_u=|y_u|v$.

Las hipercargas de cada término suman cero:
\[
  -Y_Q+Y_\phi+Y_{d_R}=0,\qquad
  -Y_Q+Y_{\widetilde\phi}+Y_{u_R}=0.
\]
Para comprobar la invariancia bajo $SU(2)$ usamos
$\sigma_2\sigma_i^*=-\sigma_i\sigma_2$. Si
$U=e^{ig\vec\alpha\cdot\vec\sigma/2}$, entonces
\[
  i\sigma_2U^*=U\,i\sigma_2,
  \qquad
  \widetilde\phi\longmapsto
  i\sigma_2(U\phi)^*=U\widetilde\phi.
\]
Así, $\widetilde\phi$ transforma en la representación fundamental,
y $\bar Q\widetilde\phi$ es un singlete de $SU(2)$.

El modelo estándar consiste de tres generaciones de leptones: $e$ (electrón),
$\mu$ (muón), $\tau$ (tauón), con interacciones y simetrías idénticas a las del
electrón. Todos vienen en versión izquierda y derecha, y tienen sus respectivos
neutrinos $\nu_e,\nu_\mu,\nu_\tau$, sólo izquierdos en el modelo estándar. Como
todos transforman igual, podemos escribir más Yukawas invariantes:
\[
  \mathcal L_{\mathrm{leptones}}
  =i\bar L_i\not\!D L_i+i\bar e_{Ri}\not\!D e_{Ri}
   -\left(y_{ij}^{(l)}\bar L_i\phi e_{Rj}+\hc\right),
\]
donde $i=1,2,3$ corresponde a $e,\mu,\tau$ respectivamente.

Para los quarks tenemos tres familias: $u$ (up) y $d$ (down), $c$ (charm) y
$s$ (strange), $t$ (top) y $b$ (bottom). De nuevo las tres transforman
idénticamente (recordemos que además de la interacción electrodébil estos
también tienen interacción de color $SU(3)$):
\[
  \mathcal L_{\mathrm{quarks}}
  =i\bar Q_i\not\!D Q_i
   +i\bar u_{Ri}\not\!D u_{Ri}
   +i\bar d_{Ri}\not\!D d_{Ri}
   -\left(
     y_{ij}^{(d)}\bar Q_i\phi d_{Rj}
     +y_{ij}^{(u)}\bar Q_i\widetilde\phi u_{Rj}
     +\hc\right).
\]

En los quarks se contraen además los índices de color, y
\[
  D_\mu=\partial_\mu-ig'YX_\mu-igJ^aW_\mu^a
        -ig_sT^AG_\mu^A,
  \qquad T^A=\lambda^{(A)}/2.
\]
Para leptones y para el Higgs, el último término está ausente.
El índice de familia $i=1,2,3$ no es un índice de color.

Reuniendo los términos y acoplando la teoría mínimamente a la gravedad,
obtenemos
\begin{align*}
  S_{\text{todo lo que sabemos}}
  =\int d^4x\sqrt{-g}\Bigg[&
    \frac{M_{\mathrm{Pl}}^2}{2}R-M_{\mathrm{Pl}}^2\Lambda
    -(D_\mu\phi)^\dagger D^\mu\phi
    -m^2\phi^\dagger\phi
    -\lambda(\phi^\dagger\phi)^2\\
  &+i\bar L_i\not\!D L_i+i\bar e_{Ri}\not\!D e_{Ri}
    +i\bar Q_i\not\!D Q_i\\
  &+i\bar u_{Ri}\not\!D u_{Ri}
    +i\bar d_{Ri}\not\!D d_{Ri}\\
  &-\Bigl(
    y_{ij}^{(l)}\bar L_i\phi e_{Rj}
    +y_{ij}^{(d)}\bar Q_i\phi d_{Rj}
    +y_{ij}^{(u)}\bar Q_i\widetilde\phi u_{Rj}
    +\hc\Bigr)\\
  &-\frac14X_{\mu\nu}X^{\mu\nu}
    -\frac14W_{\mu\nu}^aW^{a\mu\nu}
    -\frac14G_{\mu\nu}^AG^{A\mu\nu}
  \Bigg].
\end{align*}
En espacio curvo, $\not\!D=\gamma^\mu(x)D_\mu$ incluye la conexión
espinorial y las matrices gamma se construyen con una tétrada.
La normalización de $\Lambda$ coincide con la introducción.
Hemos omitido los términos topológicos, que no alteran las ecuaciones
clásicas locales para coeficientes constantes.

A esto le falta: algo que le dé masa a los neutrinos, algo que sea materia
oscura, algo que haga inflación, algo que explique la asimetría
materia--antimateria.

¡Hay mucho que no está en $S_{\text{todo lo que sabemos}}$!

\end{document}
